%==============================================================================
%  Sina Sabzevar -- Academic CV
%  Positioning: Causal discovery & effective connectivity in clinical fMRI
%
%  BUILD:  pdflatex sina_sabzevar_cv.tex   (run twice)
%  Compiles on Overleaf as-is (standard packages only).
%
%  >>> Red [FILL: ...] markers are placeholders you MUST replace before sending.
%  >>> Set \draftfalse below to preview without the red highlighting.
%==============================================================================
\documentclass[10pt,a4paper]{article}

\usepackage[T1]{fontenc}
\usepackage[utf8]{inputenc}
\usepackage{lmodern}
\usepackage[margin=0.7in,top=0.55in,bottom=0.55in]{geometry}
\usepackage{enumitem}
\usepackage{titlesec}
\usepackage{xcolor}
\usepackage{array}
\usepackage[hidelinks]{hyperref}

%------------------------------------------------------------------ colors ---
\definecolor{accent}{HTML}{1F3A5F}   % deep slate blue
\definecolor{rule}{HTML}{B8C4D4}
\definecolor{muted}{HTML}{555555}
\definecolor{todo}{HTML}{C0392B}

%--------------------------------------------------------------- draft mode ---
\newif\ifdraft
\drafttrue          % <-- \draftfalse to preview without red markers

% Breakable placeholder (must NOT be a box, or it overflows the margin)
\newcommand{\fillme}[1]{%
  \ifdraft{\color{todo}\bfseries[FILL: #1]}\else\textbf{[#1]}\fi}

%------------------------------------------------------------------ layout ---
\setlength{\parindent}{0pt}
\pagestyle{empty}
\setlength{\emergencystretch}{2em}

\titleformat{\section}
  {\color{accent}\normalsize\bfseries\scshape}
  {}{0em}{}[\vspace{-0.55em}{\color{rule}\rule{\linewidth}{0.7pt}}]
\titlespacing*{\section}{0pt}{8pt}{4pt}

\setlist[itemize]{leftmargin=1.0em,itemsep=1pt,topsep=1.5pt,parsep=0pt,
                  label=\textcolor{accent}{\textbullet}}

% Role header:  \role{Title}{Institution / advisor}{Dates}
\newcommand{\role}[3]{%
  \par\vspace{4pt}%
  \textbf{#1}\hfill{\color{muted}\small #3}\par
  \textit{\color{muted}\small #2}\par\vspace{1.5pt}}

% Compact one-line entry (teaching)
\newcommand{\teachline}[2]{%
  \par #1\hfill{\color{muted}\small #2}\par\vspace{2pt}}

% Publication group heading
\newcommand{\pubhead}[1]{\par\vspace{4pt}{\bfseries\small #1}\par\vspace{1.5pt}}

\newcounter{pub}
\newcommand{\pub}[1]{\par\refstepcounter{pub}\noindent%
  \begin{minipage}[t]{0.032\linewidth}\raggedright\small\thepub.\end{minipage}%
  \begin{minipage}[t]{0.958\linewidth}\small #1\end{minipage}\\[2.5pt]}

\newcommand{\me}{\textbf{Sabzevar, S.}}

%==============================================================================
\begin{document}

%--------------------------------------------------------------------- head ---
\begin{center}
  {\LARGE\bfseries\color{accent} Sina Sabzevar}\\[3pt]
  {\large Causal Discovery \& Effective Connectivity in Clinical fMRI}\\[6pt]
  {\small
    \href{mailto:sinasabzevar7@gmail.com}{sinasabzevar7@gmail.com}
    \;\textbar\; +98 916 746 6673
    \;\textbar\; \href{https://linkedin.com/in/sinasabzevar}{linkedin.com/in/sinasabzevar}\\[1.5pt]
    \fillme{website} \;\textbar\; \fillme{Google Scholar} \;\textbar\; \fillme{GitHub} \;\textbar\; \fillme{ORCID}
  }
\end{center}

%------------------------------------------------------------------ summary ---
\section{Research Summary}

I develop causal-inference methods for resting-state fMRI to identify discriminative brain regions
in psychiatric and neurological disorders. Across four years and two laboratories I have applied
constraint-based, score-based, and functional causal discovery (PC, GES, LiNGAM) to schizophrenia
and epilepsy cohorts, and developed generative-augmentation and graph-embedding methods to address
the small-sample, class-imbalanced regime that limits clinical neuroimaging. MSc in Electrical
Engineering, University of Tehran; ranked 4th nationally among $\sim$20,000 candidates. I currently
apply large-scale imbalanced-learning methods in industry while continuing neuroimaging research,
and seek a PhD extending causal connectivity analysis toward
\fillme{your direction, e.g.\ cross-disorder biomarker discovery}.

%---------------------------------------------------------------- education ---
\section{Education}

\role{M.Sc., Electrical Engineering (Control Systems)}{University of Tehran, Iran \ $\cdot$ \ Full Scholarship \ $\cdot$ \ GPA 16.45/20 (3.72/4.00)}{Sep 2019 -- Feb 2023}
\begin{itemize}
  \item \textbf{Thesis:} \textit{Diagnosing Discriminative Brain Regions in Schizophrenia Using
        Causal Discovery and Machine Learning.} \ Advisor: Prof.\ Babak N.\ Araabi.
\end{itemize}

\role{B.Sc., Electrical Engineering (Control Systems)}{Hamedan University of Technology, Iran \ $\cdot$ \ Full Scholarship}{Sep 2013 -- Jul 2018}

%------------------------------------------------------------- publications ---
\section{Publications}

\pubhead{Peer-Reviewed Conference Papers}
\pub{\me, Masoudnia, S., Nazemzadeh, M.R., Araabi, B.N. (2023). Enhanced Diagnosis of
  Schizophrenia Using Integrated Node Embedding and Causal Discovery on Imbalanced fMRI Data.
  \textit{IEEE Int.\ Conf.\ on Signal Processing and Intelligent Systems (ICSPIS)}. \fillme{DOI}}

\pubhead{Manuscripts Under Review}
\pub{\me, Masoudnia, S., Araabi, B.N. Rest-fMRI Data Augmentation Using Deep Generative Models to
  Improve Classification Performance. Under review at \fillme{journal}, submitted \fillme{mon.\ year}.}

\pubhead{Conference Abstracts \& Presentations}
\pub{\me, Masoudnia, S., Nazemzadeh, M.R., Araabi, B.N. (2025). Identifying Discriminative Regions
  in Epilepsy Through Effective Connectivity Analysis. \textit{Organization for Human Brain Mapping
  (OHBM) Annual Meeting}, \fillme{city}; \fillme{poster/oral}.}

\pubhead{In Preparation}
\pub{\me, \fillme{co-authors}. Continual Learning: Review, Benchmark Dataset, and Evaluation Under
  Non-Stationary Data Distributions. Target: \fillme{venue}, \fillme{2027}.}

%--------------------------------------------------------- research exper. ---
\section{Research Experience}

\role{Research Assistant --- Causal Discovery in Psychiatric Neuroimaging}{Computational Modeling \& Machine Learning Lab, University of Tehran \ $\cdot$ \ Advisor: Prof.\ B.N.\ Araabi}{2021 -- 2024}
\begin{itemize}
  \item Applied constraint-based (PC), score-based (GES), and functional (LiNGAM) causal discovery
        to resting-state fMRI from \fillme{dataset} (\fillme{n} patients / \fillme{n} controls,
        \fillme{atlas}), recovering directed connectivity that isolated \fillme{k} discriminative
        regions in schizophrenia.
  \item Integrated graph node embedding with causal-graph features, improving \fillme{metric}
        from \fillme{baseline} to \fillme{result} over correlation-based functional-connectivity
        baselines.
  \item Addressed severe class imbalance (\fillme{ratio}) with SMOTE and cost-sensitive learning,
        raising minority-class \fillme{recall/F1} by \fillme{delta}.
  \item Designed VAE- and GAN-based augmentation for rs-fMRI connectivity data, improving held-out
        \fillme{metric} by \fillme{delta} in the low-$n$ regime.
  \item Built and ran the end-to-end pipeline on HPC infrastructure (\fillme{scale}).
\end{itemize}

\role{Research Assistant --- Effective Connectivity in Epilepsy}{Research Center for Molecular \& Cellular Imaging, TUMS \ $\cdot$ \ Advisors: Prof.\ M.R.\ Nazemzadeh, Dr.\ S.\ Masoudnia}{2022 -- 2023}
\begin{itemize}
  \item Estimated effective connectivity on a clinical epilepsy cohort (\fillme{n} patients /
        \fillme{n} controls, \fillme{scanner}), identifying \fillme{k} regions distinguishing
        \fillme{subtype} from controls; presented at OHBM 2025.
  \item Developed a multimodal \fillme{modalities} pipeline for \fillme{task}, reaching
        \fillme{metric}.
  \item Collaborated directly with clinicians on cohort definition, quality control, and
        interpretation of connectivity findings.
\end{itemize}

%-------------------------------------------------------- applied / industry ---
\section{Applied Machine Learning Experience}

\role{Data Scientist \& Project Manager}{Hamrah-e Aval (MCI), Tehran}{Feb 2024 -- Present}
\begin{itemize}
  \item Built end-to-end supervised pipelines for churn prediction and subscriber classification
        over \textbf{70M+ records} --- large-scale learning under severe class imbalance, the same
        methodological regime as clinical neuroimaging at a different scale.
  \item Developed a simulation framework for counterfactual sensitivity analysis of campaigns,
        estimating intervention effects prior to deployment.
\end{itemize}

\role{Freelance Data Scientist \& Python Developer}{Remote \ $\cdot$ \ scalable data pipelines, computer vision, predictive modeling}{2020 -- 2024}
\role{Co-Founder \& Lead Instructor}{Cofe Deep --- online AI education platform; built deep-learning curriculum, delivered live lectures}{2020 -- 2021}

%------------------------------------------------------------------ teaching ---
\section{Teaching --- University of Tehran}

\teachline{\textbf{Chief Teaching Assistant} --- Machine Learning \& Deep Learning (Graduate)}{Fall 2024}
\teachline{\textbf{Teaching Assistant} --- Deep Learning (Graduate)}{Fall 2021, Fall 2022}
\teachline{\textbf{Teaching Assistant} --- Statistical Inference \& Game Theory (Graduate)}{2020}
\teachline{\textbf{Instructor} --- PyTorch \& Python hands-on workshops}{2021, 2022}

%-------------------------------------------------------------------- skills ---
\section{Technical Skills}

\begin{tabular}{@{}>{\bfseries}p{0.20\linewidth}p{0.77\linewidth}@{}}
Causal \& \newline Connectivity &
  Causal discovery (PC, GES, FCI, LiNGAM), effective connectivity, functional connectivity,
  graph \& network analysis, brain parcellation \\[3pt]
Neuroimaging &
  Resting-state fMRI preprocessing \& analysis, structural MRI,
  \fillme{name your actual toolboxes: fMRIPrep / FSL / SPM / nilearn / CONN / Tetrad} \\[3pt]
Machine Learning &
  PyTorch, scikit-learn, graph neural networks, VAEs, GANs, imbalanced learning, continual learning \\[3pt]
Computing &
  Python (advanced), MATLAB, SQL, HPC \& distributed computing, Git, \LaTeX{} \\
\end{tabular}

%-------------------------------------------------------------------- honors ---
\section{Honors \& Awards}
\begin{itemize}
  \item \textbf{Ranked 4th nationally} out of $\sim$20,000 candidates, Iranian MSc entrance
        examination (Konkour), 2019.
  \item Full Scholarship --- MSc, University of Tehran (2019--2023); BSc, Hamedan University of
        Technology (2013--2018).
\end{itemize}

%---------------------------------------------------------------- additional ---
\section{Additional Information}

\begin{tabular}{@{}>{\bfseries}p{0.20\linewidth}p{0.77\linewidth}@{}}
Languages     & English --- TOEFL iBT \fillme{retake and update} \ $\cdot$ \ Persian (native) \\[2pt]
Certification & Machine Learning Specialization, Stanford University (Coursera) \\[2pt]
Coursework    & Deep Generative Models, Advanced Deep Learning, Statistical Inference,
                System Identification, Optimal \& Non-Linear Control, Game Theory \\
\end{tabular}

%---------------------------------------------------------------- references ---
\section{References}

\begin{tabular}{@{}p{0.27\linewidth}p{0.24\linewidth}>{\raggedright\arraybackslash\small}p{0.46\linewidth}@{}}
\textbf{Prof.\ Babak N.\ Araabi}       & University of Tehran & \href{mailto:araabi@ut.ac.ir}{araabi@ut.ac.ir} \\[1.5pt]
\textbf{Dr.\ M.R.\ Nazem-Zadeh}        & Monash University    & \href{mailto:Mohamad.Nazem-Zadeh@monash.edu}{Mohamad.Nazem-Zadeh@monash.edu} \\[1.5pt]
\textbf{Dr.\ Saeed Masoudnia}          & University of Tehran & \href{mailto:masoudnia@ut.ac.ir}{masoudnia@ut.ac.ir} \\
\end{tabular}

\end{document}
